% TOP_PROBLEM: {"id": 30006677, "title": "Extended Riemann hypothesis for Dedekind zeta functions", "category_id": 1, "category": {"id": 1, "name": "number_theory", "display_name": "Number Theory", "description": "Properties of integers, prime numbers, Diophantine equations.", "slug": "number-theory", "order_index": 1, "created_at": "2026-07-31T15:26:25.670Z"}, "status": "open"}
% FIRST_POSTED: 2026-09-25
% ENTRY_KIND: statement_only
% !TeX program = lualatex
% Definition and sources only; this file is not an LLM solution attempt.
\documentclass[11pt]{article}
\usepackage[margin=25mm]{geometry}
\usepackage{fontspec}
\setmainfont{Latin Modern Roman}
\usepackage{amsmath,amssymb,mathtools}
\usepackage{unicode-math}
\setmathfont{Latin Modern Math}
\usepackage{xurl,hyperref}
\hypersetup{colorlinks=true,urlcolor=blue}
\setcounter{secnumdepth}{0}
\setlength{\parindent}{0pt}
\setlength{\parskip}{5pt}
\setlength{\emergencystretch}{3em}
\begin{document}
\section*{\texorpdfstring{Extended Riemann hypothesis for Dedekind zeta functions}{Extended Riemann hypothesis for Dedekind zeta functions}}\label{30006677}
\subsection{Definitions and mathematical statement}
For a number field $K$, let $\mathcal O_K$ be its ring of integers and $N\mathfrak a=\#(\mathcal O_K/\mathfrak a)$ the norm of a nonzero integral ideal. Define
\[
 \zeta_K(s)=\sum_{0\ne\mathfrak a\subseteq\mathcal O_K}(N\mathfrak a)^{-s}
 \qquad(\Re s>1),
\]
and continue meromorphically. The assertion is $\zeta_K(\rho)=0$ and $0<\Re\rho<1\Rightarrow\Re\rho=1/2$, for every finite extension $K/{\mathbb{Q}}$. The rational field itself is included.
\par\smallskip\textit{Formulation and concept references:} \href{https://www.aimath.org/WWN/rh/articles/html/18a/}{[S1]}.

\subsection{Short English statement}
For every number field, must the zeros of its ideal-counting zeta function inside the critical strip lie halfway across it?
\subsection{Sources}
\begin{itemize}
\item[S1]American Institute of Mathematics RH project, The Extended Riemann Hypothesis; undated article18a.
\url{https://www.aimath.org/WWN/rh/articles/html/18a/}.
\textit{Formulation source.}
\item[S2]A random polynomial with multiplicative coefficients is almost surely irreducible — Varjú and Xu.
\url{https://academic.oup.com/imrn/article/2026/16/rnag178/8761102}.
\textit{Further reference; not a proof endorsement.}
\item[S3]On the Extended, Generalized, and Grand Riemann Hypotheses: A Unified Approach Based on the General Properties of L-Functions — Weicun Zhang, v19.
\url{https://www.preprints.org/manuscript/202506.0481}.
\textit{Further reference; not a proof endorsement.}
\item[S4]Proof of generalized Riemann hypothesis for Dedekind zetas and Dirichlet L-functions — Andrzej Mądrecki.
\url{https://arxiv.org/abs/0706.0256}.
\textit{Further reference; not a proof endorsement.}
\end{itemize}
\end{document}
